\scanpage{24}
\mspage{4}
définition est de car $p>0$, et donc ils ne se prêtent pas à la détermination des nbs \uline{exacts} de p$^{\mathrm{ts}}$ fixes par formule de Lefschetz (mais seulement le nb mod $p$), donc \unsure{non plus} à une fondation satisfaisante des conjectures de Weil.

b) Le point de vue De Rham classique, \inserted{pour exprimer la coh. ``complexe'',} est trop exclusivement lié aux variétés \uline{non} singulières. Rien n'empêche évidemment de considérer aussi le complexe $\underline\Omega_X^\bullet$ et $\mathbb H^*(X,\Omega_X^\bullet)$ si $X$ est quelconque. Mais $\overline{\mathrm m}$ si $X$ est une courbe, affine ou complète, sur $\mathbb C$,
\deleted{singulière}
\inserted{p.ex. la droite affine ou projective \deleted{\unsure{qui a un point double}} un p$^{\mathrm t}$ de \unsure{rebr.} ordinaire,}
\deleted{(\unsure{disons unibranche si on y tient})}, on n'obtient déjà pas pour \deleted{\unsure{la}} $\dim H_{\mathrm{DR}}^1(X)$ le premier nb de Betti comme il faudrait.
On peut décider qu'on ne s'intéresse qu'aux variétés algébriques non singulières — mais l'expérience apprend que $\overline{\mathrm m}$ pour étudier celles-ci, on ne peut éviter l'introduction des variétés singulières, p.ex. à titre de fibres intervenant dans des \deleted{\unsure{fibrations}} morphismes
\margin{23}
